\VerdierPhysicalPageAnchor{0115}%
\begingroup
\small
\setlength{\abovedisplayskip}{4pt plus 2pt minus 2pt}%
\setlength{\belowdisplayskip}{4pt plus 2pt minus 2pt}%

\VerdierSemanticAnchor{ch2-1-2-3-proof-f1-f2}
Suppose, to fix ideas, that $f_1\colon X\longrightarrow X'$ and
$f_2\colon Y\longrightarrow Y'$ are isomorphisms. For every object $W$ of
$\mathcal D$, we have a commutative diagram:

\VerdierSemanticAnchor{ch2-1-2-3-five-lemma-diagram}
\begin{equation*}
  \resizebox{0.98\linewidth}{!}{%
    $\begin{tikzcd}[ampersand replacement=\&,column sep=2.15em,row sep=3em]
      \cdots \arrow[r]
        \& \operatorname{Hom}_{\mathcal D}(W,X) \arrow[r] \arrow[d,"{\operatorname{Hom}_{\mathcal D}(W,f_1)}"']
        \& \operatorname{Hom}_{\mathcal D}(W,Y) \arrow[r] \arrow[d,"{\operatorname{Hom}_{\mathcal D}(W,f_2)}"']
        \& \operatorname{Hom}_{\mathcal D}(W,Z) \arrow[r] \arrow[d,"{\operatorname{Hom}_{\mathcal D}(W,f_3)}"']
        \& \operatorname{Hom}_{\mathcal D}(W,X[1]) \arrow[r] \arrow[d,"{\operatorname{Hom}_{\mathcal D}(W,f_1[1])}"']
        \& \cdots \\
      \cdots \arrow[r]
        \& \operatorname{Hom}_{\mathcal D}(W,X') \arrow[r]
        \& \operatorname{Hom}_{\mathcal D}(W,Y') \arrow[r]
        \& \operatorname{Hom}_{\mathcal D}(W,Z') \arrow[r]
        \& \operatorname{Hom}_{\mathcal D}(W,X'[1]) \arrow[r]
        \& \cdots
    \end{tikzcd}$%
  }\,.
\end{equation*}

\VerdierSemanticAnchor{ch2-1-2-3-five-lemma}
The rows are exact. It follows, by the five lemma, that
$\operatorname{Hom}_{\mathcal D}(W,f_3)$ is an isomorphism for every $W$,
and hence that $f_3$ is an isomorphism. The other two assertions of the
corollary follow from this one by using axiom TRII.

\VerdierSemanticAnchor{ch2-1-2-4}
\medskip
\noindent\textbf{Corollary 1.2.4.}\quad
Let
$X\xrightarrow{u}Y\xrightarrow{v}Z\xrightarrow{w}X[1]$ and
$X\xrightarrow{u}Y\xrightarrow{v'}Z'\xrightarrow{w'}X[1]$
be two distinguished triangles. There exists an isomorphism
$f\colon Z\xrightarrow{\sim}Z'$ such that the following diagram is
commutative:

\VerdierSemanticAnchor{ch2-1-2-4-1}
\begin{equation}
  \begin{tikzcd}[column sep=2.8em,row sep=1.7em]
    && Z \arrow[dr,"w"] \arrow[dd,"f"',"\sim"] & \\
    X \arrow[r,"u"] & Y \arrow[ur,"v"] \arrow[dr,"v'"'] && X[1] \\
    && Z' \arrow[ur,"w'"'] &
  \end{tikzcd}
  \tag{1.2.4.1}
\end{equation}

\VerdierSemanticAnchor{ch2-1-2-4-proof}
Indeed, by axiom TRIII, there exists a morphism
$f\colon Z\longrightarrow Z'$ such that diagram (1.2.4.1) is commutative.
By (1.2.3), this morphism is an isomorphism.

\VerdierSemanticAnchor{ch2-1-2-4-essential-uniqueness}
Corollary (1.2.4) therefore shows that distinguished triangles containing a
given morphism are essentially unique: they are determined up to
isomorphism. It should be noted, however, that this isomorphism is not
unique. We shall make this point precise in (1.2.12).

\VerdierSemanticAnchor{ch2-1-2-5}
\medskip
\noindent\textbf{Corollary 1.2.5.}\quad
The direct sum of two distinguished triangles is a distinguished triangle.

\VerdierSemanticAnchor{ch2-1-2-5-proof-start}
Let:

\VerdierSemanticAnchor{ch2-1-2-5-input-triangles}
\[
  \begin{gathered}
    X\xrightarrow{u}Y\xrightarrow{v}Z\xrightarrow{w}X[1],\\
    X'\xrightarrow{u'}Y'\xrightarrow{v'}Z'\xrightarrow{w'}X'[1]
  \end{gathered}
\]
\endgroup
